% ============================================================
%  Informe 1 Práctica 2 — SOII
%  Compilar: lualatex informe1_practica2.tex  (dos pasadas)
% ============================================================
\documentclass[a4paper, 12pt]{article}

% ── Codificación y fuentes ──────────────────────────────────
\usepackage[T1]{fontenc}
\usepackage[utf8]{inputenc}

% ── Idioma ──────────────────────────────────────────────────
\usepackage[spanish, es-nolists, es-noindentfirst]{babel}

% ── Geometría ───────────────────────────────────────────────
\usepackage[
  a4paper,
  top=2.5cm, bottom=2.5cm,
  left=2.8cm, right=2.8cm
]{geometry}

% ── Colores ─────────────────────────────────────────────────
\usepackage{xcolor}
\definecolor{azul}     {HTML}{2E5FA3}
\definecolor{azulosc}  {HTML}{1A3A6B}
\definecolor{grisclaro}{HTML}{F4F4F4}
\definecolor{grisborde}{HTML}{CCCCCC}
\definecolor{rojocod}  {HTML}{A31515}

% ── Encabezados y pies ──────────────────────────────────────
\usepackage{fancyhdr}
\pagestyle{fancy}
\fancyhf{}
\fancyhead[L]{\small\color{azul}\textbf{SOII --- Práctica 2}}
\fancyhead[R]{\small\color{gray}Informe 1: Carreras críticas}
\fancyfoot[C]{\small\thepage}
\renewcommand{\headrulewidth}{0.4pt}

% ── Títulos de sección ──────────────────────────────────────
\usepackage{titlesec}
\titleformat{\section}
  {\Large\bfseries\color{azul}}
  {\thesection.}{0.6em}{}
  [\vspace{-0.3em}\textcolor{azul}{\rule{\linewidth}{0.8pt}}\vspace{0.2em}]

\titleformat{\subsection}
  {\large\bfseries\color{azulosc}}
  {\thesubsection.}{0.5em}{}

\titleformat{\subsubsection}
  {\normalsize\bfseries}
  {\thesubsubsection.}{0.4em}{}

% ── Listas ──────────────────────────────────────────────────
\usepackage{enumitem}
\setlist[itemize]{leftmargin=1.5em, itemsep=2pt, topsep=4pt}

% ── Tablas ──────────────────────────────────────────────────
\usepackage{booktabs}
\usepackage{array}
\usepackage{colortbl}
\newcolumntype{L}[1]{>{\raggedright\arraybackslash}p{#1}}
\newcolumntype{C}[1]{>{\centering\arraybackslash}p{#1}}

% ── Código fuente ───────────────────────────────────────────
\usepackage{listings}
\lstdefinestyle{cstyle}{
  language=C,
  basicstyle=\ttfamily\small,
  keywordstyle=\bfseries\color{azul},
  commentstyle=\itshape\color{gray},
  stringstyle=\color{rojocod},
  numberstyle=\tiny\color{gray},
  numbers=left,
  numbersep=8pt,
  stepnumber=1,
  backgroundcolor=\color{grisclaro},
  frame=single,
  framesep=4pt,
  rulecolor=\color{grisborde},
  breaklines=true,
  tabsize=4,
  showstringspaces=false,
  xleftmargin=1.8em,
  xrightmargin=0.5em,
  captionpos=b,
  literate=
    {á}{{\'a}}1 {é}{{\'e}}1 {í}{{\'i}}1 {ó}{{\'o}}1 {ú}{{\'u}}1
    {Á}{{\'A}}1 {É}{{\'E}}1 {Í}{{\'I}}1 {Ó}{{\'O}}1 {Ú}{{\'U}}1
    {ñ}{{\~n}}1 {Ñ}{{\~N}}1,
}
\lstdefinestyle{shellstyle}{
  basicstyle=\ttfamily\small,
  backgroundcolor=\color{grisclaro},
  frame=single,
  framesep=4pt,
  rulecolor=\color{grisborde},
  breaklines=true,
  numbers=none,
  xleftmargin=1em,
  xrightmargin=0.5em,
}
\lstset{style=cstyle}

% ── Cajas resaltadas ────────────────────────────────────────
\usepackage{tcolorbox}
\tcbuselibrary{skins, breakable}

\newtcolorbox{infobox}[1][]{
  colback=azul!7, colframe=azul, fonttitle=\bfseries,
  title=#1, left=6pt, right=6pt, top=4pt, bottom=4pt,
  breakable
}
\newtcolorbox{warnbox}[1][]{
  colback=orange!8, colframe=orange!70!black, fonttitle=\bfseries,
  title=#1, left=6pt, right=6pt, top=4pt, bottom=4pt,
  breakable
}

% ── Hiperenlaces ────────────────────────────────────────────
\usepackage{hyperref}
\hypersetup{
  colorlinks=true, linkcolor=azul, urlcolor=azul,
  pdftitle={Práctica 2 SOII --- Informe 1},
}

% ── Miscelánea ──────────────────────────────────────────────
\usepackage{microtype}
\usepackage{parskip}
\setlength{\parskip}{6pt}

% ── Macro inline-code ───────────────────────────────────────
\newcommand{\ic}[1]{\texttt{#1}}

% ============================================================
\begin{document}

% ── Portada ─────────────────────────────────────────────────
\begin{titlepage}
  \centering
  \vspace*{3cm}
  {\Huge\bfseries\color{azul} Sistemas Operativos II\par}
  \vspace{0.6cm}
  {\LARGE Práctica 2 --- Informe 1\par}
  \vspace{0.4cm}
  {\large\color{gray} Carreras críticas en el productor-consumidor\par}
  \vspace{0.5cm}
  \textcolor{azul}{\rule{0.7\linewidth}{1.2pt}}
  \vspace{2cm}

  \begin{tabular}{rl}
    \textbf{Asignatura:} & Sistemas Operativos II \\[4pt]
    \textbf{Titulación:} & Grado en Ingeniería Informática \\[4pt]
    \textbf{Curso:}      & 2024/25 \\[4pt]
    \textbf{Fecha:}      & \today \\
  \end{tabular}

  \vfill
  {\small Universidad de Santiago de Compostela}
\end{titlepage}

\tableofcontents
\newpage

% ============================================================
\section{Informe 1 --- Carreras críticas en el productor-consumidor}
% ============================================================

\subsection{Descripción del problema}

El problema del productor-consumidor es un problema clásico de concurrencia
en el que dos procesos comparten un buffer de tamaño fijo (\ic{N~=~10}).
El productor escribe caracteres en el buffer y el consumidor los retira.
Cuando no existe ningún mecanismo de sincronización, ambos procesos pueden
acceder simultáneamente a las variables compartidas \ic{count} y \ic{top},
produciendo \emph{carreras críticas}.

\subsection{Cómo ejecutar los programas}

\subsubsection*{Compilación}
\begin{lstlisting}[style=shellstyle]
gcc -o prod prod.c
gcc -o cons cons.c
\end{lstlisting}

\subsubsection*{Ejecución (dos terminales distintas)}
\begin{lstlisting}[style=shellstyle]
# Terminal 1:
./prod
# Terminal 2 (arrancar poco despues del productor):
./cons
\end{lstlisting}

El fichero de texto \ic{input.txt} debe encontrarse en el mismo directorio
de trabajo. El productor debe arrancarse primero para que cree el fichero de
memoria compartida \ic{/tmp/shm\_buf}.

\subsection{Diseño y funcionamiento}

\subsubsection*{Memoria compartida}

La memoria compartida se implementa mediante \ic{mmap()} sobre el fichero
\ic{/tmp/shm\_buf}. La estructura compartida contiene:

\begin{itemize}
  \item \ic{buf[N]}: buffer de tipo \ic{char} que funciona como pila LIFO.
  \item \ic{top}: índice de la cima de la pila (\ic{top = 0} significa pila
        vacía).
  \item \ic{count}: número de elementos actualmente en el buffer.
\end{itemize}

\subsubsection*{Funciones principales}

\begin{itemize}
  \item \ic{produce\_item()}: lee un carácter del fichero \ic{input.txt}
        e incrementa el contador de la vocal correspondiente.
  \item \ic{insert\_item(item)}: inserta \ic{item} en \ic{buf[top]} y
        actualiza \ic{top} y \ic{count}. \textbf{Región crítica.}
  \item \ic{remove\_item()}: decrementa \ic{top}, lee el carácter de la
        cima, deja un guión \ic{'-'} en su lugar y decrementa \ic{count}.
        \textbf{Región crítica.}
  \item \ic{consume\_item(item)}: contabiliza las vocales del carácter
        retirado.
\end{itemize}

\subsubsection*{Espera activa}

Al no disponer de semáforos, la sincronización se realiza mediante
\emph{busy-waiting}: el productor comprueba \ic{count == N} en un bucle
y el consumidor comprueba \ic{count == 0}, durmiendo un segundo en cada
iteración de espera.

\subsection{Región crítica y cómo provocar carreras}

La región crítica comprende el acceso combinado a \ic{count} y \ic{top}.
El código incluye líneas \ic{sleep()} comentadas marcadas con
\ic{[SLEEP-RC]} que permiten ampliar artificialmente la duración de dicha
región:

\begin{lstlisting}[style=cstyle,
  caption={Fragmento de \texttt{insert\_item()} en \texttt{prod.c}}]
void insert_item(char item)
{
    /* ... espera activa si buffer lleno ... */

    /* sleep(SLEEP_RC); */  /* [SLEEP-RC] amplifica la carrera */

    shm->buf[shm->top] = item;
    shm->top++;
    shm->count++;
}
\end{lstlisting}

\begin{warnbox}[Cómo activar las carreras críticas]
  Modificar la constante \texttt{SLEEP\_RC} en \texttt{prod.c} y \texttt{cons.c}.
  Al introducir un \texttt{sleep()} dentro de la región crítica, un proceso se 
  congela antes de terminar de actualizar \texttt{top} y \texttt{count}, permitiendo 
  al otro proceso entrar y operar sobre datos desfasados.
\end{warnbox}

\subsection{Resultados observados y patrones de carrera}

Al modificar los valores de retardo dentro de la región crítica (\ic{SLEEP\_RC}), se observaron patrones de carrera muy específicos, fuertemente condicionados por las velocidades relativas de los procesos, el límite del buffer (\ic{N=10}) y el número total de iteraciones (\ic{MAX\_ITER=80}):

\begin{itemize}
  \item \textbf{Productor lento (\ic{sleep=1}) y Consumidor rápido (\ic{sleep=0}):}\\
        Se observan carreras críticas de forma constante. El productor pasa un segundo entero ``congelado'' a mitad de su actualización de variables, lo que deja una ventana de vulnerabilidad enorme. El consumidor, al no tener retardo, entra repetidas veces a la región crítica y colisiona continuamente con la operación incompleta del productor.
        
  \item \textbf{Consumidor lento (\ic{sleep=1}) y Productor rápido (\ic{sleep=0}):}\\
        Las carreras se presentan de forma asimétrica. El productor \textbf{no} sufre carreras en sus primeras 10 iteraciones porque el buffer inicial está vacío y las inserta rápidamente antes de que el consumidor despierte de su primer letargo. Inversamente, el consumidor \textbf{no} sufre carreras en sus últimas 10 iteraciones porque el productor ya ha terminado sus 80 ciclos totales y abandonado la memoria compartida, dejando al consumidor operando en solitario.
        
  \item \textbf{Ambos rápidos (\ic{sleep=0} en los dos):}\\
        Las carreras críticas ocurren casi exclusivamente en las iteraciones múltiples de 10. Al ejecutarse sin retardos artificiales, los procesos colisionan cuando alcanzan los límites físicos del buffer (0 o 10 elementos). Estos topes obligan a un proceso a entrar en \emph{busy-waiting} frente a la región crítica. En el momento exacto en que el otro proceso altera el contador, ambos se abalanzan sobre la memoria compartida simultáneamente, garantizando la colisión exactamente cada 10 iteraciones.
\end{itemize}

El recuento final de vocales (cuando el código se lograba ejecutar sin corrupción fatal) fue idéntico en ambos procesos:

\vspace{0.5em}
\begin{center}
\begin{tabular}{C{3cm} C{3cm} C{3cm}}
  \toprule
  \rowcolor{azul!80}
  \color{white}\textbf{Vocal} &
  \color{white}\textbf{Producidas} &
  \color{white}\textbf{Consumidas} \\
  \midrule
  a\,/\,A & 8 & 8 \\
  e\,/\,E & 6 & 6 \\
  i\,/\,I & 8 & 8 \\
  o\,/\,O & 5 & 5 \\
  u\,/\,U & 3 & 3 \\
  \bottomrule
\end{tabular}
\end{center}

\subsection{Conclusiones}

\begin{itemize}
  \item Sin sincronización, las carreras críticas son posibles siempre
        que la región crítica no sea atómica desde el punto de vista del
        planificador del sistema operativo.
  \item Los límites físicos del buffer actúan como cuellos de botella que fuerzan la sincronización temporal de los procesos, multiplicando la probabilidad de colisiones sistemáticas.
  \item El valor de \ic{count} puede superar \ic{N} o hacerse negativo
        si ambos procesos lo modifican concurrentemente.
  \item El índice \ic{top} puede apuntar a una posición errónea,
        provocando sobreescritura o lectura de datos corruptos en el
        buffer.
  \item La espera activa (\emph{busy-waiting}) desperdicia CPU y no
        garantiza exclusión mutua: es una solución incorrecta e
        ineficiente.
\end{itemize}

\end{document}